\documentclass[ ../main.tex]{subfiles}
\providecommand{\mainx}{..}
\begin{document}
\section{Random approximate values}
\label{sec:approx_values}

A Bernoulli set approximates set membership: the indicator function $\SetIndicator{\Set{A}} \colon \Set{U} \to \BitSet$ may err on any element with quantifiable probability.
A Bernoulli map generalizes this to arbitrary functions $\Fun{f} \colon \Set{X} \to \Set{Y}$, where the output for each input may independently differ from the true output.
In both cases, the fundamental object is the same: a \emph{value} that may differ from some latent (true) value with quantifiable probability.
This section formalizes that object.

\subsection{Definition}

\begin{definition}[Random approximate value]
\label{def:approx_value}
Let $\Set{T}$ be a type (a finite set of values).
A \emph{random approximate value} of type $\Set{T}$ is a random variable $\tilde{x}$ taking values in $\Set{T}$, associated with a latent value $x^* \in \Set{T}$, such that
\begin{equation}
\label{eq:approx_value_pmf}
\Prob{\tilde{x} = t \Given x^* = s} = M_{s,t}
\end{equation}
for all $s, t \in \Set{T}$, where $M$ is a $\Card{\Set{T}} \times \Card{\Set{T}}$ row-stochastic matrix called the \emph{confusion matrix} (or \emph{channel matrix}) of the approximation.
We write $\tilde{x} \in \AT{\Set{T}}$ to denote that $\tilde{x}$ is an approximate value of type $\Set{T}$.
\end{definition}

The confusion matrix $M$ encodes the complete error model.
Its diagonal entries $M_{s,s}$ are the probabilities of correct observation; its off-diagonal entries $M_{s,t}$ for $s \neq t$ are the probabilities of specific confusions.
A general $\Card{\Set{T}} \times \Card{\Set{T}}$ row-stochastic matrix has $\Card{\Set{T}}(\Card{\Set{T}} - 1)$ free parameters.
In the Bernoulli framework, structural assumptions (independence, partition into blocks) reduce this count, yielding parsimonious models parameterized by a small number of error rates.

\begin{remark}[Connection to Bernoulli models]
When $\Card{\Set{T}} = 2$, the confusion matrix is a $2 \times 2$ channel.
If $\Set{T} = \BitSet = \{0, 1\}$, the off-diagonal entries are the false positive rate $\fprate = M_{0,1}$ and the false negative rate $\fnrate = M_{1,0}$.
This is precisely the second-order Bernoulli Boolean from the companion paper on Bernoulli models.
The notation $\AT{\Set{T}}[\fprate][\fnrate]$ makes the rates explicit when $\Set{T}$ is binary.
\end{remark}

\begin{remark}[$\BitSet$ vs.\ $\Bool$]
\label{rem:bitset_bool}
Throughout this paper, $\BitSet = \{0, 1\}$ and $\Bool = \{\True, \False\}$ are isomorphic two-element sets used in different contexts.
We write $\BitSet$ in numerical and information-theoretic settings (confusion matrices, error rates, channel models) and $\Bool$ in logical and type-theoretic settings (Boolean algebra, logic gates, type constructors).
No distinction in the underlying error model is intended; the choice of notation reflects the surrounding context.
\end{remark}

\subsection{The primary mechanism: approximate maps}

The standard way an approximate value arises is through function application.
An \emph{approximate map} $\APFun{f} \colon \Set{X} \to \Set{Y}$ approximates a latent function $\Fun{f} \colon \Set{X} \to \Set{Y}$ by independently introducing per-input errors.
When $\APFun{f}$ is applied to an exact input $x \in \Set{X}$, the output $\APFun{f}(x)$ is an approximate value of type $\Set{Y}$:
\begin{equation}
\APFun{f}(x) \in \AT{\Set{Y}}\,.
\end{equation}
The confusion matrix for $\APFun{f}(x)$ is the per-input channel $M_x$, a $\Card{\Set{Y}} \times \Card{\Set{Y}}$ row-stochastic matrix.
This is the exponential-type viewpoint: an approximate map is an approximate value of the function type $\Set{X} \to \Set{Y}$, with the joint confusion matrix factoring as a Kronecker product of per-input channels (\cref{sec:exponential}).

\subsection{Monadic structure}

Approximate values compose.
If $\Fun{g} \colon \Set{A} \to \Set{B}$ is an exact function and $\tilde{a} \in \AT{\Set{A}}$ is an approximate value, then $\Fun{g}(\tilde{a})$ is an approximate value of type $\Set{B}$:
\begin{equation}
\label{eq:monadic_lift}
\Fun{g}(\tilde{a}) \in \AT{\Set{B}}\,.
\end{equation}
More generally, if both $\Fun{g}$ and the input are approximate, the composition yields a higher-order approximate value.
This gives $\AT{(\cdot)}$ the structure of a \emph{monad}: the operation of applying a function to an approximate value (the monadic \emph{bind}) produces an approximate value whose error rates are determined by the composition of the respective confusion matrices.

We make this precise in \cref{sec:error_prop}, where we derive the output confusion matrix for function application on approximate inputs.

\subsection{Worked example: the approximate bit}
\label{sec:approx_bitset}

The simplest non-trivial type is $\BitSet = \{0, 1\}$.
An approximate value $\tilde{b} \in \AT{\BitSet}[\fprate][\fnrate]$ is parameterized by two error rates:
\begin{itemize}
\item $\fprate = \Prob{\tilde{b} = 1 \Given b^* = 0}$ (false positive rate), and
\item $\fnrate = \Prob{\tilde{b} = 0 \Given b^* = 1}$ (false negative rate).
\end{itemize}
The full conditional probability mass function is given by the $2 \times 2$ confusion matrix:

\begin{center}
\begin{tabular}{@{}lcc@{}}
\toprule
& observe $0$ & observe $1$ \\
\midrule
latent $0$ & $1 - \fprate$ & $\fprate$ \\
latent $1$ & $\fnrate$ & $1 - \fnrate$ \\
\bottomrule
\end{tabular}
\end{center}

Equivalently, the conditional PMF is
\begin{equation}
\label{eq:bitset_pmf_0}
\Fun{p}(b \Given 0) =
\begin{cases}
1 - \fprate & \text{if } b = 0\,, \\
\fprate     & \text{if } b = 1\,,
\end{cases}
\end{equation}
and
\begin{equation}
\label{eq:bitset_pmf_1}
\Fun{p}(b \Given 1) =
\begin{cases}
\fnrate     & \text{if } b = 0\,, \\
1 - \fnrate & \text{if } b = 1\,.
\end{cases}
\end{equation}

This is the Bernoulli Boolean: the building block for all approximate types with $\Card{\Set{T}} \geq 2$.
When the approximate bit arises from a Bernoulli set membership query, $\fprate$ and $\fnrate$ are the set's false positive and false negative rates.

\subsection{Lifting functions through approximate values}
\label{sec:lifting_preview}

Suppose $\Fun{g} \colon \BitSet \to \BitSet$ is an exact function and $\tilde{b} \in \AT{\BitSet}[\fprate][\fnrate]$ is an approximate bit.
Then $\Fun{g}(\tilde{b})$ is an approximate bit whose conditional PMF depends on $\Fun{g}$.

Given latent value $b^* = 0$, the input $\tilde{b}$ realizes $0$ with probability $1 - \fprate$ and $1$ with probability $\fprate$, so
\begin{equation}
\label{eq:lift_cond_0}
\Fun{p}_{\Fun{g}(\tilde{b})}(c \Given 0) =
\begin{cases}
1 - \fprate & \text{if } c = \Fun{g}(0)\,, \\
\fprate     & \text{if } c = \Fun{g}(1)\,.
\end{cases}
\end{equation}

Given latent value $b^* = 1$, the input $\tilde{b}$ realizes $0$ with probability $\fnrate$ and $1$ with probability $1 - \fnrate$, so
\begin{equation}
\label{eq:lift_cond_1}
\Fun{p}_{\Fun{g}(\tilde{b})}(c \Given 1) =
\begin{cases}
\fnrate     & \text{if } c = \Fun{g}(0)\,, \\
1 - \fnrate & \text{if } c = \Fun{g}(1)\,.
\end{cases}
\end{equation}

For the identity function $\Fun{g} = \IdFun$, the output rates equal the input rates.
For the constant functions ($\Fun{g} \equiv 0$ or $\Fun{g} \equiv 1$), the output is exact regardless of the input error.
For the negation function $\Fun{g} = \NegateFn$, the roles of $\fprate$ and $\fnrate$ swap.
These four cases exhaust $\BitSet \to \BitSet$ and illustrate the general pattern: the output error rates are determined by the composition of the input confusion matrix with the function's action on the codomain.

\end{document}
